\documentclass[12pt]{article}
\usepackage{anysize}
\usepackage{xltxtra}
\usepackage{xunicode}
\usepackage{graphicx}
\usepackage{color}
\usepackage{tikz}
\usepackage{minted}
\usepackage{listings}
\usepackage{enumitem}
\usepackage{framed}
\usepackage{relsize}
\usepackage{float}
\usepackage{algpseudocode}
\usepackage[hyphens]{url}

\usepackage{color,hyperref}

\usepackage{fancyhdr}
\usepackage{lastpage}
\usepackage{fontspec}

\linespread{1.5}

\definecolor{darkblue}{rgb}{0.0,0.0,0.3}
\hypersetup{colorlinks,breaklinks,
        linkcolor=darkblue,urlcolor=darkblue,
            anchorcolor=black,citecolor=black}

\def\name{Grigorios Lyras}
\def\IF{\textbf{if}}
\def\THEN{\textbf{then}}
\def\SAYS{\textbf{says}}
\def\SPEAKSFOR{$\Rightarrow$}
\def\FOR{\textbf{for}}
\def\AND{\textbf{and}}

\setlength{\headheight}{15.2pt}

\pagestyle{fancy}{ %
    \fancyhf{} % remove everything
    \renewcommand{\headrulewidth}{0pt} % remove lines as well
    \renewcommand{\footrulewidth}{0.5pt}
    \lfoot{\name}
    \rfoot{Page \thepage\ of \pageref{LastPage}}

\renewcommand{\labelenumi}{(\alph{enumi})}
\defaultfontfeatures{Ligatures=TeX} % To support LaTeX quoting style
\setmainfont[Mapping=TeX-text]{CMU Serif}


\begin{document}

\begin{titlepage}

\begin{tikzpicture}[remember picture, overlay]
      \node [anchor=north east, inner sep=0pt]  at (current page.north east)
           {\includegraphics[height=3cm]{title/ucl_logo}};
       \end{tikzpicture}
\begin{center}
\begin{LARGE}\textbf{
University College London\\}\end{LARGE}
\vspace{1cm}
\begin{Large}Simulating Path Planning for Data\end{Large}\\
\begin{Large}Muling in Wireless Sensor Networks\end{Large}\\
\begin{Large}using Unmanned Aerial Vehicles\end{Large}\\
\vspace{2cm}
\begin{LARGE}
MSc Networked Computer Systems\\
NCS Group Project Individual Report\\
\end{LARGE}
\large{Academic Year 2013-2014}\\
\vfill
\end{center}
This report is submitted as part requirement for the MSc in Networked Computer
Systems at University College London. It is substantially the result of our own
work except where explicitly indicated in the text. The report may be freely
copied and distributed provided the source is explicitly acknowledged.\\
\begin{center}
\begin{tabular}{l r}
\Large{Grigorios Lyras}&
\large{SID: 13047325}\\
\end{tabular}\\
\vspace{1cm}
\large\today\\
\end{center}
\end{titlepage}



\tableofcontents{}

\pagebreak

\section{Introduction}
The aim of this report is the discussion of the Group Project that I took part
in as part of the NCS team, comprised of myself and Seng Jea Lee, Johann Mifsud,
Minmin Zhang. In this report I will present my own role in the progress of
the project. The analysis will be broken into three main sections. In objectives
and achievements I will present the objectives we set through the duration of
the project as well as the achievements and my contribution to them. Moving to
the results of the project we will discuss the output of the project as well as
the quality of the outcome. In the third section, I will present my view on the
project management techniques and tools that we used along with my experience on
them. The interaction with the rest of the team members and the task
distribution will also be reported. In the final section I will present my own
contribution to the project.

Wireless Sensor Networks are systems of sensors distributed over an area in
order to collect useful information. These are low cost devices with limitations
in their capabilities. One of these limitations is the communication range. As a
result it is understood that these networks are not always connected. A concept
to overcome such difficulties is the use of data mules that visit these devices
and collet the data from them. These data mules are networked devices that are
able to move from location to location, pick up the data, store them
temporarily and then deliver them back to the base. An interesting area is the
exploration of data muling with the use of Unmanned Aerial Vehicles (UAVs).
UAVs have unique mobility abilities which make them a very suitable solution
for data muling.

Since the suitability of UAVs for the task is established we move to the
implementation. One of the factors that increase the risk significantly in this
researh area is the cost of the UAVs as well as the dangers they impose to
the nearby environment. This danger dictates that the behaviour and performane of
UAVs needs to be tested thoroughly and the high cost makes field testing
unsuitable for this researh area.

The initial idea for this project was the exploration of path planning
algorithms for UAVs when they are used to collect data from distributed wireless
sensor networks. Such networks appear more and more often in the industry mainly
because of their lower cost
The focus of this project has been the exploration of path planning algorithms.
More specifically
Simulation of such algorithms

\section{Objectives and Achievements}
One of the most challenging aspects of this project has been it's wide scope. It
is in effect a multi disciplinary project.


your own understanding of the project objectives and main achievements;

\section{Project Results}
\subsection{Radio Modeling}
\subsubsection{ZigBee}
\subsubsection{WiFi}
\subsection{Simulator Plugin}
\subsection{Path Planning}
a critical assessment of the project results (e.g., how do you judge the
quality of the final result, what were the important technical decisions
this project, what worked well, what decisions you would change in the light
of experience, etc.);

\section{Project Management}
a critical assessment of the project management techniques used (e.g.,
whether you followed standard software engineering processes, and if so then
how well they worked or didn't work);

a discussion and critical assessment of how well the team and each of its
members worked from your point-of-view (e.g., whether decisions were
adequately agreed, specified and documented, or had to be revised later,
whether the work was appropriately split among team members);

a description and critical assessment of your own workplan and contribution
to the project (e.g., what aspects of the project did you contribute most
to; whether you were over-worked or under-worked; whether you as an
individual could have achieved more with different approaches and/or
organisation).

In the Individual report, you may make reference to the submitted Group
documents for further detail, but the report should stand on its own.

Note that the Individual report is only made available to the Examiners and will
not be seen by any other students.


\end{document}
